\documentclass[book.tex]{subfiles}

\begin{document}

\chapter{Green's Identities}

\label{chap:greens_identities}
\noindent\rule{11cm}{0.4pt} \\
\textbf{Prerequisites:}  \autoref{chap:calculus}, \autoref{chap:vectors} \\
\textbf{Difficulty Level:} *** \\
\noindent\rule{11cm}{0.4pt}
\vspace{5mm}



Green's identities are a set of three vector calculus identities used widely in physics.

\section{Green's First Identity}
Let $U \subset \mathbb{R}^n$ be a region of space with boundary $\partial U$. Let $\varphi$ and $\psi$ be scalar functions on $U$ with $\varphi$ twice differentiable and $\psi$ differentiable. Let $\vec{n}$ be an outward pointing normal vector. Then
\begin{align}
\int_U (\psi \nabla^2 \varphi + \nabla \psi \cdot \nabla \varphi) dV &= \oint_{\partial U} \psi (\nabla \varphi \cdot \vec{n}) dS \\
&= \oint_{\partial U} \psi \nabla \varphi \cdot d\vec{S}
\end{align}
where we have
\begin{align}
d\vec{S} &= \vec{n} dS
\end{align}
is the oriented surface element. This rule is a higher-dimensional analog of the technique of integration by parts.
%\textcolor{red}{TODO: Explain derivation, give an example}

\section{Green's Second Identity}

Let $\varphi$ and $\psi$ be twice continuously differentiable functions. Then we have 

\begin{align}
\int_U \left [ \psi \nabla^2  \varphi - \varphi \nabla^2  \psi   \right] dV = \oint_{\partial U} \left ( \psi \frac{\partial \varphi}{\partial \vec{n}} - \varphi \frac{\partial \psi}{\partial \vec{n}} \right ) dS
\end{align}

$\frac{\partial \varphi}{\partial \vec{n}}$ is the directional derivative of $\varphi$ in the normal direction.

\begin{align}
\frac{\partial \varphi}{\partial \vec{n}} &= \nabla \varphi \cdot \vec{n}
\end{align}


\section{Green's Third Identity}

Let $\psi$ be a function that is twice continuously differentiable functions on $U$. Let $G$ be the Green's function of the Laplacian $\nabla^2$. Then we have that

\begin{align}
\int_U [G(y, \eta) \nabla^2 \psi(y)] dV_y - \psi(\eta) &= \oint_{\partial U} \left [ G(y, \eta)\frac{\partial \psi}{\partial \vec{n}} - \psi(y) \frac{\partial G(y, \eta)}{\partial \vec{n}} \right ] dS_y
\end{align}

This formula proves useful when solving boundary value problems in electrostatics.

\section{The Divergence Theorem}

The divergence theorem provides a broader theoretical justification for the facts considered in this section.


\end{document}

